\documentclass[11pt]{article}

\usepackage[margin=1in]{geometry}
\usepackage{amsmath,amssymb,amsthm,mathtools}
\usepackage{enumitem}
\usepackage{booktabs}
\usepackage{tabularx}
\usepackage{array}
\usepackage{microtype}
\usepackage[hidelinks]{hyperref}
\usepackage[nameinlink,noabbrev]{cleveref}
\usepackage{xurl}

\newtheorem{definition}{Definition}
\newtheorem{proposition}{Proposition}
\newtheorem{remark}{Remark}
\newcolumntype{Y}{>{\raggedright\arraybackslash}X}

\newcommand{\defeq}{\stackrel{\mathrm{def}}{=}}
\newcommand{\MaxL}{\mathcal{L}}

\title{Conditional Per-Step Budget Certificates\\
\large A Receipt-Grade Interface for Adaptive Horizon Claims (Series Synthesis)}
\author{}
\date{Unified archive note v0.05 (2026-02-28; archive v123)}

\begin{document}
\maketitle

\begin{abstract}
The archive publishes privacy budgets as compact receipt line items with digest-bound replay hooks.
For repeated use (retries, fallbacks, and interactive protocols), additive horizon claims do \emph{not} require independence,
but they \emph{do} require a conditional per-step guarantee.
This note defines a small, auditable object shape---a \emph{conditional per-step budget certificate}---that is sufficient to justify
horizon-$n$ receipts under adaptive interaction.
Other papers should cite this note instead of re-explaining what ``replayable per-step conditional bounds'' means.
\end{abstract}

\paragraph{Scope.}
This note is an interface: it defines what must be committed to (and replayable) for Regime~II horizon accounting.
It is not a new privacy theorem; the sequential-composition proof is owned by Anonymity~A.\cite{series:oddsinflation}

\paragraph{Imports.}
Receipt line-item fields (including \texttt{replay\_hook}) are owned by Synthesis~12.\cite{series:receiptitems}
Horizon-regime rules (I/II/III) are owned by Synthesis~19.\cite{series:accountingrulebook}
Plan/state drift semantics for digest-bound hooks are owned by Synthesis~18.\cite{series:planstateequiv}
The mechanical replay checklist is owned by Synthesis~20.\cite{series:auditorreplay}

\paragraph{Later-terse-paper citation posture.}
When a later short paper only needs to answer \emph{which conditional per-step certificate shape justifies this adaptive horizon claim?}, it should stop at the minimal public certificate card below plus the tiny horizon drill.
It should widen only when the live issue becomes the sequential-composition theorem root (Anonymity~A), the receipt line-item tuple itself (Synthesis~12), regime classification (Synthesis~19), plan/state drift semantics (Synthesis~18), the mechanical replay recipe (Synthesis~20), or spectral / runtime-specific evidence translations (Anonymity~D and Operational~B).\cite{series:oddsinflation,series:receiptitems,series:accountingrulebook,series:planstateequiv,series:auditorreplay,series:spectral2budget,series:operationalB}

\begin{table}[t]
\centering
\small
\begin{tabularx}{0.97\linewidth}{@{}p{0.30\linewidth}Y@{}}
\toprule
\textbf{Public row} & \textbf{Pinned value}\\
\midrule
Receipt scope & Regime~II adaptive horizon claim over $T=30$ steps; the public receipt line item will publish only the total horizon budget plus a replay hook \\
Conditioning interface & \texttt{cond\_iface\_id} commits to the deterministic summary $S_{t-1}=g_t(H_{t-1})$ used at every step; changing that interface mints a new certificate family rather than a silent patch \\
Per-step cap & Each step carries a declared conditional cap \texttt{bound\_bits = 0.03} for the realized output under the realized public conditioning value \\
Evidence hook & \texttt{evidence\_id} points to either a static certificate, a finite conditioning table, or a replayable program committed by \texttt{plan\_spec\_id} \\
Replay binding & The receipt's \texttt{replay\_hook} binds \texttt{plan\_id}, \texttt{plan\_spec\_id}, and \texttt{cond\_iface\_id}, so an auditor can recompute the per-step checks from public logs and requested artifacts \\
Drift discipline & If the conditioning interface or declared state surface changes, Synthesis~18's digest-binding / state-equivalence discipline decides whether the old certificate still applies or a new claim must be issued \\
Owner routing & Anonymity~A owns the composition theorem; Synthesis~12 owns the line-item tuple; Synthesis~20 owns replay mechanics; spectral / runtime-specific translations remain in Anonymity~D and Operational~B \\
\bottomrule
\end{tabularx}
\caption{A minimal public conditional per-step budget certificate card. This is the smallest public answer later terse papers can quote directly when the live issue is whether an adaptive horizon receipt really rests on replayable conditional step caps rather than on an independence assumption.}
\label{tab:minimal-cps-card}
\end{table}

This card is intentionally non-decorative: it pins one conditioning interface, one per-step cap family, one evidence hook discipline, and one replay binding, without silently re-deriving the theorem root or the broader release stack.

\paragraph{One tiny adaptive-horizon drill.}
If a receipt claims $T=30$ adaptive steps with a constant declared cap $b_t=0.03$ bits per step, then the horizon budget justified by Proposition~\ref{prop:suff} is simply
\[
\sum_{t=1}^{30} b_t = 30\cdot 0.03 = 0.90\ \text{bits}.
\]
A later terse paper may therefore say ``the horizon-30 receipt relied on replayable conditional step caps of $0.03$ bits each, for a total declared budget of $0.90$ bits'' and move on, unless the live issue is the theorem proof, the line-item schema, or the replay mechanics.

\section{Why this certificate exists}
Anonymity~A gives a sequential-composition theorem: if each step satisfies a conditional one-vs-mixture bound given the public history,
then the total odds-inflation (equivalently, pointwise MaxL in bits) is bounded by the sum of per-step budgets.\cite{series:oddsinflation}
Receipts want to publish horizon claims (``$Q\times$'') without assuming independence; this is Regime~II in Synthesis~19.\cite{series:accountingrulebook}
To keep papers short and referee-grade, we isolate the \emph{auditable commitment} that makes the theorem applicable.

\section{Conditional per-step budget certificates}
We write $I$ for the hidden identity (or secret) and $O_t$ for the public output at step $t$.
Let $H_{t-1}$ denote the public history (all public outputs so far, plus the declared schedule/plan information).

\begin{definition}[Conditioning interface]\label{def:conditioning}
A conditional certificate must specify a public \emph{conditioning interface}:
a deterministic projection
\[
S_{t-1}\ \defeq\ g_t(H_{t-1})
\]
together with a digest identifier \texttt{cond\_iface\_id} that binds the code/specification of $g_t$.
Intuitively, $S_{t-1}$ is the public summary of history that the per-step guarantee conditions on.
\end{definition}

\begin{definition}[Conditional per-step MaxL certificate]\label{def:cps}
Fix a step $t$.
A \emph{conditional per-step MaxL certificate} for $t$ consists of:
\begin{enumerate}[leftmargin=*]
\item a bound value $b_t\ge 0$ (in \emph{bits}, i.e.\ $\log_2$ units),
\item a conditioning interface \texttt{cond\_iface\_id} (Definition~\ref{def:conditioning}),
\item a digest-bound \emph{evidence hook} \texttt{evidence\_id} that enables an auditor to check the inequality below for the realized conditioning value $S_{t-1}$.
\end{enumerate}
The certificate asserts that, almost surely over the public history, the realized output satisfies
\begin{equation}\label{eq:condmaxl}
\max_i \log_2 \frac{\Pr[O_t=o\mid I=i,\,S_{t-1}]}{\Pr[O_t=o\mid S_{t-1}]}\ \le\ b_t
\qquad\text{for the realized }o=O_t.
\end{equation}
The evidence hook may be implemented by (i) a static dual/cut certificate (when the kernel is history-invariant),
(ii) a finite table over conditioning states, or (iii) a replayable program committed by \texttt{plan\_spec\_id} that deterministically computes the checked inequality from public logs and requested artifacts.
\end{definition}

\begin{remark}[A minimal record shape for tool output]
When tooling emits per-step certificates (for use inside a \texttt{plan\_spec}), a compact record format is:
\[
\{\ \texttt{step\_selector},\ \texttt{cond\_iface\_id},\ \texttt{bound\_bits},\ \texttt{evidence\_id}\ \}.
\]
Here \texttt{step\_selector} names which step(s) the bound applies to (a single $t$, a range, or a predicate on public schedule time), and \texttt{evidence\_id} points to either (i) a static dual/cut certificate, (ii) a finite conditioning table, or (iii) a log-audit object.
The equalization-audit lemma in Anonymity~B is an example of a log-audit evidence object that can populate \texttt{evidence\_id} for a discrete projection.\cite{series:profiling}
\end{remark}



\begin{remark}[Where this lives in a receipt]
Receipts do \emph{not} need new fields.
A line item that claims a horizon-$Q$ total under Regime~II must include a \texttt{replay\_hook}
(\texttt{plan\_id}, \texttt{plan\_spec\_id}) whose plan specification commits to:
(i) the conditioning interface \texttt{cond\_iface\_id},
(ii) the per-step bounds $b_t$ (or a deterministic rule that yields them), and
(iii) the evidence hooks \texttt{evidence\_id} needed for replay.
This is an instance of the general replay discipline in Synthesis~20.\cite{series:auditorreplay}
\end{remark}

\begin{proposition}[Receipt-facing sufficiency]\label{prop:suff}
If a receipt's replay hook provides conditional per-step MaxL certificates (Definition~\ref{def:cps}) for steps $t=1,\dots,T$,
then the horizon-$T$ transcript satisfies a total MaxL/odds-inflation budget of at most $\sum_{t=1}^T b_t$ bits.
\end{proposition}

\begin{remark}[Stateful deployments]
For anonymous DHTs, routing tables, caches, congestion signals, and scores can make $S_{t-1}$ large.
Regime~III in Synthesis~19 treats this by enlarging the witness family via a digest-bound \texttt{state\_decl\_id}.\cite{series:accountingrulebook,series:planstateequiv}
A Regime~II certificate remains usable in the stateful setting as long as the conditioning interface $g_t$ is defined over the declared state surfaces
(and the replay plan obtains the relevant state readouts from the transparency bundle).
\end{remark}

\section{Research warranted (kept out of other papers)}
Two recurring barriers to small certificates are (a) huge conditioning spaces and (b) optional stopping / drift monitoring.
For (a), a promising direction is to certify contraction of the conditioning channel (e.g.\ via a spectral gap or a mixing-rate certificate)
so one can bound worst-case conditional leakage by a small state proxy.\cite{series:spectral}
For (b), Operational~B's anytime-valid evidence suggests ``drift certificates'' that remain safe under optional stopping.\cite{series:operationalB}
We keep those research threads out of mechanism papers; this note is the single citable landing place.

\subsection{Spectral contraction as a compact evidence object (pointer)}
\label{sec:spectral-cert}

One concrete way to keep Regime~II certificates small is to certify \emph{contraction} of a delegation kernel rather than enumerate a huge conditioning space.
For random-walk delegation, the archive's published spectral anonymity note provides sharp distinguishing bounds and conservative certificates (e.g.\ via a second-singular-value bound).\cite{series:spectral}

This note does \emph{not} re-derive spectral translations.
The canonical ``spectral witness $\Rightarrow$ per-step MaxL cap (bits)'' recipe is Anonymity~D.\cite{series:spectral2budget}
Mechanism papers that want to export a Regime~II per-step certificate from a delegation walk should cite Anonymity~D for the numeric translation and cite this note for the receipt-side certificate shape.

\begin{thebibliography}{99}\small

\bibitem{series:oddsinflation}
Anonymous.
\newblock Odds-Inflation Anonymity: Posterior Odds Budgets, Composition, and Operational Limits.
\newblock Anonymity series Paper~A, The-One-True-Archive (2026).

\bibitem{series:receiptitems}
Anonymous.
\newblock Receipt Line Items for Anonymity Claims: A Minimal Tuple Schema for Published Budgets.
\newblock Synthesis~12, The-One-True-Archive (2026).

\bibitem{series:accountingrulebook}
Anonymous.
\newblock Receipt Accounting and Horizon Composition for Anonymous-DHT Budgets: A Rulebook for Repeats, Retries, and State.
\newblock Synthesis~19, The-One-True-Archive (2026).

\bibitem{series:planstateequiv}
Anonymous.
\newblock Plan and State Equivalence: Digest Binding, Drift, and Replay Semantics.
\newblock Synthesis~18, The-One-True-Archive (2026).

\bibitem{series:auditorreplay}
Anonymous.
\newblock Auditor Replay for Anonymous-DHT Receipt Line Items: A Mechanical Checklist for Digest-Bound Claims.
\newblock Synthesis~20, The-One-True-Archive (2026).

\bibitem{series:spectral}
Anonymous.
\newblock Spectral Anonymity and Optimal Distinguishing Bounds for Random-Walk Delegation in (Possibly Directed) Overlays.
\newblock Published note (2026-01-23), The-One-True-Archive (2026).

\bibitem{series:spectral2budget}
Anonymous.
\newblock Spectral Delegation for Anonymous DHT Lookups: From mixing certificates to auditable per-step budgets.
\newblock Anonymity series Paper~D, The-One-True-Archive (2026).

\bibitem{series:operationalB}
Anonymous.
\newblock Auditable Trace Privacy with Abort: Disclosure receipts for side-channel resistant transport.
\newblock Operational series Paper~B, The-One-True-Archive (2026).
\bibitem{series:profiling}
Anonymous.
\newblock Anonymous DHT Profiling and the Equalization Contract.
\newblock Anonymity series Paper~B, The-One-True-Archive (2026).
\end{thebibliography}

\end{document}
